\documentclass[aspectratio=169]{beamer}
\usepackage[utf8]{inputenc}
\usepackage{xeCJK}
\usepackage{graphicx}
\usepackage{booktabs}
\usepackage{tikz}
\usepackage{pgfplots}
\usepackage{xcolor}
\usepackage{colortbl}
\usepackage{longtable}
\usepackage{array}

% 設定中文字體
\setCJKmainfont{Noto Sans CJK TC}

% 主題設定
\usetheme{Madrid}
\usecolortheme{seahorse}

% 標題頁資訊
\title{心臟內科可信賴專業活動EPAs}
\subtitle{Cardiology Entrustable Professional Activities}
\author{內科部心臟科}
\date{\today}
\institute{醫學中心}

% 自訂顏色 - 沿用原文件配色
\definecolor{emergency_red}{RGB}{186,24,27}
\definecolor{emergency_blue}{RGB}{0,114,188}
\definecolor{emergency_gray}{RGB}{120,120,120}
\definecolor{highlight_yellow}{RGB}{255,250,205}

% 調整beamer顏色主題
\setbeamercolor{structure}{fg=emergency_blue}
\setbeamercolor{frametitle}{fg=emergency_red}
\setbeamercolor{block title}{bg=emergency_blue!20,fg=emergency_blue}
\setbeamercolor{block body}{bg=emergency_blue!5}
\setbeamercolor{block title alerted}{bg=emergency_red!20,fg=emergency_red}
\setbeamercolor{block body alerted}{bg=emergency_red!5}

\begin{document}

% 標題頁
\begin{frame}
\titlepage
\end{frame}

% 大綱
\begin{frame}{今日議程}
\tableofcontents
\end{frame}

\section{心臟內科EPAs概述}

\begin{frame}{什麼是心臟內科EPAs？}
\begin{block}{Entrustable Professional Activities (EPA)}
\textbf{可信賴專業活動} - 心臟內科醫師在臨床實務中能夠獨立勝任的專業工作單元
\end{block}

\vspace{0.5cm}

\begin{columns}[c]
\column{0.5\textwidth}
\textbf{涵蓋範圍：}
\begin{itemize}
\item 心血管疾病診斷
\item 急慢性心臟病處理
\item 心臟檢查技術
\item 病患照護協調
\end{itemize}

\column{0.5\textwidth}
\textbf{核心特色：}
\begin{itemize}
\item 以能力為導向
\item 漸進式發展
\item 真實情境評量
\item 個別化學習
\end{itemize}
\end{columns}
\end{frame}

\begin{frame}{委託等級架構}
\begin{center}
\begin{tikzpicture}
% 定義顏色
\definecolor{level1}{RGB}{255,182,193}
\definecolor{level2}{RGB}{255,218,185}
\definecolor{level3}{RGB}{255,255,224}
\definecolor{level4}{RGB}{144,238,144}
\definecolor{level5}{RGB}{173,216,230}

% 繪製發展階梯
\fill[level1] (0,0) rectangle (2,1);
\node at (1,0.5) {\scriptsize Level 1};
\node at (1,-0.5) {\scriptsize Novice};

\fill[level2] (2,0) rectangle (4,1);
\node at (3,0.5) {\scriptsize Level 2};
\node at (3,-0.5) {\scriptsize Advanced};
\node at (3,-0.7) {\scriptsize Beginner};

\fill[level3] (4,0) rectangle (6,1);
\node at (5,0.5) {\scriptsize Level 3};
\node at (5,-0.5) {\scriptsize Competent};

\fill[level4] (6,0) rectangle (8,1);
\node at (7,0.5) {\scriptsize Level 4};
\node at (7,-0.5) {\scriptsize Proficient};

\fill[level5] (8,0) rectangle (10,1);
\node at (9,0.5) {\scriptsize Level 5};
\node at (9,-0.5) {\scriptsize Expert};

% 邊框
\draw[thick] (0,0) rectangle (2,1);
\draw[thick] (2,0) rectangle (4,1);
\draw[thick] (4,0) rectangle (6,1);
\draw[thick] (6,0) rectangle (8,1);
\draw[thick] (8,0) rectangle (10,1);

% 箭頭
\draw[->, thick] (0,-1) -- (10,-1);
\node at (5,-1.5) {能力發展進程};

% 標示重點
\node[above] at (2,1.2) {\scriptsize 需督導};
\node[above] at (6,1.2) {\scriptsize 獨立作業};
\node[above] at (9,1.2) {\scriptsize 指導他人};
\end{tikzpicture}
\end{center}

\vspace{0.5cm}
\begin{itemize}
\item \textbf{Level 1-2}：在督導下學習與執行
\item \textbf{Level 3}：能獨立勝任標準作業
\item \textbf{Level 4-5}：能處理複雜情況並指導他人
\end{itemize}
\end{frame}

\section{10大EPAs架構}

\begin{frame}{心臟內科10大EPAs概覽}
\footnotesize
\begin{table}
\centering
\begin{tabular}{|p{1cm}|p{4cm}|p{6cm}|}
\hline
\rowcolor{emergency_blue!20}
\textbf{EPA} & \textbf{專業活動} & \textbf{核心能力領域} \\
\hline
EPA 1 & 病史詢問與風險評估 & 症狀評估、危險因子、功能分級 \\
\hline
EPA 2 & 心血管理學檢查 & 聽診技巧、血管檢查、水腫評估 \\
\hline
EPA 3 & 心電圖判讀分析 & 心律判斷、ST-T變化、傳導系統 \\
\hline
EPA 4 & 心臟超音波操作 & 影像取得、功能測量、異常識別 \\
\hline
EPA 5 & 急性冠心症處理 & 快速診斷、風險分層、緊急治療 \\
\hline
\end{tabular}
\end{table}
\end{frame}

\begin{frame}{心臟內科10大EPAs概覽 (續)}
\footnotesize
\begin{table}
\centering
\begin{tabular}{|p{1cm}|p{4cm}|p{6cm}|}
\hline
\rowcolor{emergency_blue!20}
\textbf{EPA} & \textbf{專業活動} & \textbf{核心能力領域} \\
\hline
EPA 6 & 心律不整診療 & 心律識別、血流評估、治療選擇 \\
\hline
EPA 7 & 心衰竭管理 & 診斷分類、藥物調整、整合照護 \\
\hline
EPA 8 & 高血壓管理 & 診斷確認、風險評估、生活介入 \\
\hline
EPA 9 & 心導管判讀 & 解剖識別、狹窄評估、治療規劃 \\
\hline
EPA 10 & 照護協調溝通 & 團隊合作、病患衛教、困難溝通 \\
\hline
\end{tabular}
\end{table}
\end{frame}

\section{重點EPAs詳述}

\begin{frame}{EPA 1: 心臟科病史詢問與風險評估}
\begin{alertblock}{核心描述}
能夠獨立進行心血管疾病相關的完整病史詢問，包括症狀評估、危險因子識別、家族史調查，並能初步評估心血管風險。
\end{alertblock}

\begin{columns}[c]
\column{0.5\textwidth}
\textbf{關鍵技能：}
\begin{itemize}
\item 胸痛特質詳細詢問
\item NYHA功能分級
\item 三高危險因子評估
\item 家族史系統調查
\end{itemize}

\column{0.5\textwidth}
\textbf{里程碑重點：}
\begin{itemize}
\item Level 2: 系統性病史詢問
\item Level 3: 風險分層評估
\item Level 4: 複雜情況處理
\item Level 5: 指導他人技巧
\end{itemize}
\end{columns}
\end{frame}

\begin{frame}{EPA 2: 心臟血管系統理學檢查}
\begin{alertblock}{核心描述}
能夠熟練且準確地進行心臟血管系統的完整理學檢查，包括心臟聽診、血管檢查、水腫評估，並能正確解讀檢查發現。
\end{alertblock}

\begin{columns}[c]
\column{0.5\textwidth}
\textbf{關鍵技能：}
\begin{itemize}
\item 四區心臟聽診
\item 心雜音識別分類
\item 頸靜脈壓評估
\item 周邊血管搏動
\end{itemize}

\column{0.5\textwidth}
\textbf{里程碑重點：}
\begin{itemize}
\item Level 2: 常見異常識別
\item Level 3: 複雜雜音區分
\item Level 4: 細微徵象偵測
\item Level 5: 教學示範能力
\end{itemize}
\end{columns}
\end{frame}

\begin{frame}{EPA 3: 心電圖判讀與分析}
\begin{alertblock}{核心描述}
能夠獨立判讀12導程心電圖，識別正常與異常變化，包括心律不整、ST-T變化、傳導異常等，並提出臨床相關性分析。
\end{alertblock}

\begin{columns}[c]
\column{0.5\textwidth}
\textbf{關鍵技能：}
\begin{itemize}
\item 基本間期測量
\item 心律不整分類
\item ST段異常判讀
\item 傳導阻滯識別
\end{itemize}

\column{0.5\textwidth}
\textbf{里程碑重點：}
\begin{itemize}
\item Level 2: 明顯異常識別
\item Level 3: 臨床關聯整合
\item Level 4: 複雜心律判讀
\item Level 5: 罕見變化識別
\end{itemize}
\end{columns}
\end{frame}

\begin{frame}{EPA 5: 急性冠心症診斷與初步處理}
\begin{alertblock}{核心描述}
能夠快速識別急性冠心症候群，進行適當的診斷性檢查，評估風險分層，並執行初步治療措施，包括藥物治療與轉介決策。
\end{alertblock}

\begin{columns}[c]
\column{0.5\textwidth}
\textbf{關鍵技能：}
\begin{itemize}
\item 典型/非典型症狀
\item TIMI/GRACE評分
\item 抗血小板治療
\item PCI適應症評估
\end{itemize}

\column{0.5\textwidth}
\textbf{里程碑重點：}
\begin{itemize}
\item Level 2: 標準流程執行
\item Level 3: 獨立診療決策
\item Level 4: 複雜個案處理
\item Level 5: 專業諮詢提供
\end{itemize}
\end{columns}
\end{frame}

\begin{frame}{EPA 7: 心衰竭診斷與治療}
\begin{alertblock}{核心描述}
能夠診斷急性與慢性心衰竭，區分收縮性與舒張性功能異常，制定個別化治療計畫，包括藥物調整與非藥物治療。
\end{alertblock}

\begin{columns}[c]
\column{0.5\textwidth}
\textbf{關鍵技能：}
\begin{itemize}
\item HFrEF/HFpEF區分
\item 利尿劑調整策略
\item ACE-I/ARB/ARNI選擇
\item 心臟復健規劃
\end{itemize}

\column{0.5\textwidth}
\textbf{里程碑重點：}
\begin{itemize}
\item Level 2: 基本藥物調整
\item Level 3: 複雜組合管理
\item Level 4: 難治性處理
\item Level 5: 個人化策略
\end{itemize}
\end{columns}
\end{frame}

\section{評量方法}

\begin{frame}{EPA評量工具對應}
\begin{table}
\centering
\footnotesize
\begin{tabular}{|l|l|l|}
\hline
\rowcolor{emergency_blue!20}
\textbf{EPA類型} & \textbf{主要工具} & \textbf{輔助方法} \\
\hline
技能操作類 & DOPS & 直接觀察評估 \\
(EPA 2, 4) & \scriptsize Direct Observation & 同儕回饋 \\
\hline
臨床推理類 & Mini-CEX & 病例討論 \\
(EPA 1,3,5,6,7,8,9) & \scriptsize Mini-Clinical Exercise & 模擬情境 \\
\hline
溝通協調類 & MSF & 360度回饋 \\
(EPA 10) & \scriptsize Multi-Source Feedback & 病患家屬評量 \\
\hline
\end{tabular}
\end{table}

\vspace{0.5cm}

\begin{block}{評量頻率建議}
\begin{itemize}
\item \textbf{每月重點評量}：2-3個EPAs
\item \textbf{季度全面檢討}：所有EPAs進展
\item \textbf{年度能力認證}：達標確認與規劃
\end{itemize}
\end{block}
\end{frame}

\begin{frame}{月度評量主題安排}
\begin{center}
\begin{tikzpicture}
% 繪製時間軸
\draw[thick, emergency_blue] (0,0) -- (12,0);

% 月份標記
\foreach \x/\month in {0/第1月, 3/第2月, 6/第3月, 9/第4月, 12/第5月} {
  \draw[emergency_blue] (\x,-0.2) -- (\x,0.2);
  \node[below, emergency_blue] at (\x,-0.3) {\scriptsize \month};
}

% EPAs安排
\node[above, align=center] at (1.5,0.8) {\scriptsize EPA 1,2,3\\基礎評估};
\node[above, align=center] at (4.5,0.8) {\scriptsize EPA 4,5,6\\診斷急症};
\node[above, align=center] at (7.5,0.8) {\scriptsize EPA 7,8,10\\慢病溝通};
\node[above, align=center] at (10.5,0.8) {\scriptsize EPA 9\\整合評量};

% 連接線
\foreach \x in {1.5, 4.5, 7.5, 10.5} {
  \draw[dashed, emergency_gray] (\x,0.2) -- (\x,0.6);
}
\end{tikzpicture}
\end{center}

\vspace{1cm}

\begin{alertblock}{重點提醒}
每月選擇特定EPAs進行深度評量，確保學習者能循序漸進建立專業能力。
\end{alertblock}
\end{frame}

\section{實施策略}

\begin{frame}{CCC評量架構}
\begin{block}{Case-based Clinical Competency 特色}
\begin{itemize}
\item \textbf{真實情境評量}：在實際臨床工作中評估
\item \textbf{多次觀察}：持續追蹤能力發展
\item \textbf{漸進委託}：依能力逐步增加責任
\item \textbf{360度回饋}：多方人員提供回饋
\end{itemize}
\end{block}

\vspace{0.5cm}

\begin{columns}[c]
\column{0.5\textwidth}
\textbf{評量重點：}
\begin{itemize}
\item 安全性第一
\item 能力漸進發展
\item 學習從錯誤改進
\item 個別化指導
\end{itemize}

\column{0.5\textwidth}
\textbf{成功關鍵：}
\begin{itemize}
\item 明確期望設定
\item 及時具體回饋
\item 學習資源支持
\item 持續改善機制
\end{itemize}
\end{columns}
\end{frame}

\begin{frame}{實施時程規劃}
\begin{center}
\begin{tikzpicture}
% 時間軸
\draw[thick] (0,0) -- (12,0);

% 階段標記
\foreach \x/\phase in {0/準備期, 3/試行期, 6/實施期, 9/檢討期, 12/推廣期} {
  \draw (\x,-0.1) -- (\x,0.1);
  \node[below] at (\x,-0.3) {\scriptsize \phase};
}

% 活動內容
\node[above, align=center] at (1.5,1) {\scriptsize 教師訓練\\EPA設計};
\node[above, align=center] at (4.5,1) {\scriptsize 小規模試行\\系統調整};
\node[above, align=center] at (7.5,1) {\scriptsize 全面實施\\持續評量};
\node[above, align=center] at (10.5,1) {\scriptsize 成效分析\\經驗分享};

% 里程碑
\foreach \x/\milestone in {3/設計完成, 6/試行成功, 9/實施穩定, 12/持續改善} {
  \fill[emergency_red] (\x,0) circle (0.1);
  \node[rotate=45, anchor=west, emergency_red] at (\x,0.2) {\tiny \milestone};
}
\end{tikzpicture}
\end{center}

\vspace{1cm}

\begin{block}{關鍵里程碑}
\begin{itemize}
\item \textbf{3個月}：完成EPA設計與教師訓練
\item \textbf{6個月}：試行評估並調整系統
\item \textbf{9個月}：全面實施並穩定運作
\item \textbf{12個月}：檢討成效並持續改善
\end{itemize}
\end{block}
\end{frame}

\begin{frame}{學習資源與支持}
\begin{columns}[c]
\column{0.5\textwidth}
\textbf{教學資源：}
\begin{itemize}
\item 標準化病患
\item 模擬教學設備
\item 線上學習平台
\item 心導管影像庫
\item 心電圖題庫
\end{itemize}

\textbf{評量工具：}
\begin{itemize}
\item 數位評量平台
\item 行動裝置App
\item 即時回饋系統
\item 學習歷程追蹤
\end{itemize}

\column{0.5\textwidth}
\textbf{支持系統：}
\begin{itemize}
\item Mentor制度
\item 同儕學習小組
\item 跨科際交流
\item 國際經驗分享
\end{itemize}

\textbf{品質保證：}
\begin{itemize}
\item 評量者訓練
\item 標準化流程
\item 定期校正會議
\item 外部品質稽核
\end{itemize}
\end{columns}
\end{frame}

\section{預期效益}

\begin{frame}{實施EPAs的預期效益}
\begin{block}{教育品質提升}
\begin{itemize}
\item 學習目標更明確具體
\item 評量方式更多元實用
\item 個別化學習支持
\item 能力認證更客觀
\end{itemize}
\end{block}

\begin{block}{臨床照護改善}
\begin{itemize}
\item 醫師能力標準化
\item 病人安全保障
\item 醫療品質提升
\item 團隊協作強化
\end{itemize}
\end{block}

\begin{alertblock}{長期影響}
建立以能力為導向的醫學教育文化，培養更優秀的心臟內科專科醫師
\end{alertblock}
\end{frame}

\begin{frame}{成功實施的關鍵因素}
\begin{center}
\begin{tikzpicture}
% 中心EPA圓圈
\node[circle, fill=emergency_blue, text=white, font=\Large\bfseries, minimum size=2.5cm] at (0,0) {EPAs\\成功實施};

% 四個關鍵因素
\node[rectangle, fill=emergency_red!20, text=black, minimum width=2.8cm, minimum height=1.2cm] at (-3.5,2) {
\begin{tabular}{c}
\textbf{領導支持} \\
高層承諾 \\
資源投入
\end{tabular}
};

\node[rectangle, fill=emergency_red!20, text=black, minimum width=2.8cm, minimum height=1.2cm] at (3.5,2) {
\begin{tabular}{c}
\textbf{教師準備} \\
培訓充分 \\
共識建立
\end{tabular}
};

\node[rectangle, fill=emergency_red!20, text=black, minimum width=2.8cm, minimum height=1.2cm] at (-3.5,-2) {
\begin{tabular}{c}
\textbf{系統整合} \\
流程標準 \\
技術支持
\end{tabular}
};

\node[rectangle, fill=emergency_red!20, text=black, minimum width=2.8cm, minimum height=1.2cm] at (3.5,-2) {
\begin{tabular}{c}
\textbf{持續改善} \\
定期檢討 \\
動態調整
\end{tabular}
};

% 連接線
\foreach \angle in {45, 135, 225, 315} {
  \draw[thick, emergency_blue] (0,0) -- (\angle:2.8cm);
}
\end{tikzpicture}
\end{center}
\end{frame}

\section{討論與展望}

\begin{frame}{討論議題}
\begin{enumerate}
\item 現有教學資源是否足夠支持EPAs實施？
\item 評量者訓練的具體規劃與時程安排？
\item 與現行住院醫師評量制度如何整合？
\item 學習者接受度與適應性評估？
\item 跨科EPAs的協調與統一標準？
\item 國際認證與接軌的考量？
\end{enumerate}

\vspace{1cm}
\begin{center}
\Large{\textcolor{emergency_blue}{\textbf{歡迎提問與討論}}}
\end{center}
\end{frame}

\begin{frame}{未來展望}
\begin{alertblock}{短期目標 (1年內)}
\begin{itemize}
\item 完成10個心臟內科EPAs的標準化設計
\item 建立完整的評量工具與流程
\item 訓練足夠的合格評量者
\item 建置數位化學習與評量平台
\end{itemize}
\end{alertblock}

\begin{block}{中期目標 (2-3年)}
\begin{itemize}
\item 擴展到其他內科次專科
\item 建立跨院校合作網絡
\item 發展創新教學方法
\item 進行教學成效研究
\end{itemize}
\end{block}

\begin{exampleblock}{長期願景 (5年以上)}
\begin{itemize}
\item 成為國內醫學教育標竿
\item 與國際標準接軌認證
\item 推廣至其他醫學專科
\item 建立持續改善機制
\end{itemize}
\end{exampleblock}
\end{frame}

\begin{frame}{結語}
\begin{center}
\Large{\textcolor{emergency_red}{\textbf{心臟內科EPAs}}}
\end{center}

\vspace{0.5cm}

\begin{center}
\large{不只是評量工具，更是教育革新的起點}
\end{center}

\vspace{1cm}

\begin{itemize}
\item 從知識傳授到能力培養
\item 從標準化到個別化學習
\item 從單向教學到互動發展
\item 從孤立評估到整合成長
\end{itemize}

\vspace{1cm}

\begin{center}
\textbf{\textcolor{emergency_blue}{讓我們攜手建立更優質的心臟內科教育！}}
\end{center}
\end{frame}

\begin{frame}[plain]
\begin{center}
\Huge{\textcolor{emergency_blue}{\textbf{謝謝聆聽}}}

\vspace{1cm}

\Large{\textcolor{emergency_gray}{Questions \& Discussion}}

\vspace{1cm}

\textcolor{emergency_red}{\textbf{內科部心臟科}}\\
\textcolor{emergency_blue}{Cardiology Division, Department of Internal Medicine}
\end{center}
\end{frame}

\end{document}